\documentclass[a4paper]{article}
\usepackage{parskip}
\usepackage{lipsum}

\def\nterm {Spring}
\def\nyear {2026}
\def\ncourse {Probability -- Problems}

\newcommand{\convd}{\Rightarrow}
\newcommand{\convw}{\xrightarrow{w}}
\newcommand{\convp}{\xrightarrow{p}}

\input{../header}
\author{\nauthor\vspace{5pt}\\
Carnegie Mellon University}

\newcommand{\TODO}{\textcolor{red}{\textbf{*** TO-DO ***}}}

\begin{document}
\maketitle

\tableofcontents

\setcounter{section}{1}

\section{Laws of large numbers}

\setcounter{subsection}{1}

\subsection{Weak laws of large numbers}

\begin{ex}[Durrett 2.2.1]
  Let $X_1, X_2, \dots$ be uncorrelated random variables with
  finite expectations and
  \[
    i^{-1} \var X_i \to 0 \text{ as $i \to \infty$}.
  \]
  Let $S_n = X_1 + \dots + X_n$ and $\mu_n = E S_n / n$. Show that
  $S_n / n - \mu_n \to 0$ in probability as $n \to \infty$.
\end{ex}


\begin{proof}
  WLOG assume $E X_n = 0$ for all $n$. Now we just
  need to show that $S_n / n \to 0$ in probability.
  Fix $\delta > 0$, we next show
  \[
    P \left( \abs{\frac{S_n}{n}} > 2 \delta \right)
    \to 0
  \]

  Let $\epsilon > 0$ be arbitrary.
  Since $i^{-1} \var X_i \to 0$ as $i \to \infty$,
  there exists $N$ such that $i \geq N$ implies
  $i^{-1} \var X_i < \epsilon$.
  Let $n \geq N$, we have
  \begin{align*}
    P \left( \abs{\frac{S_n}{n}} > 2 \delta \right)
     & = P \left( \frac{1}{n} \abs{\sumi^n X_i} > 2 \delta \right) \\
     & \leq P \left( \frac{1}{n} \sumi^N \abs{X_i} +
    \frac{1}{n} \abs{\sum_{i=N+1}^{n} X_i} > 2 \delta \right)      \\
     & \leq P \left( \frac{1}{n} \sumi^N \abs{X_i} >
    \delta \right)
    + P \left( \frac{1}{n} \abs{\sum_{i=N + 1}^n X_i} >
    \delta \right)
  \end{align*}
  For the first term, Chebyshev's inequality gives
  \begin{align*}
    P \left( \frac{1}{n} \sumi^N \abs{X_i} > \delta \right)
     & \leq \frac{1}{n \delta} \sumi^N E \abs{X_i} \to 0
  \end{align*}
  as $n \to \infty$.
  For the second term, Chebyshev's inequality gives
  \begin{align*}
    P \left( \frac{1}{n} \abs{\sum_{i=N+1}^n X_i} > \delta \right)
    \leq \frac{1}{\delta^2} \var \left( \frac{1}{n} \sum_{i=N + 1}^n X_i \right)
    = \frac{1}{n^2 \delta^2} \var \left( \sum_{i=N+1}^n X_i \right).
  \end{align*}
  Since $X_i$'s are uncorrelated, it follows that
  \begin{align*}
    P \left( \frac{1}{n} \abs{\sum_{i=N+1}^n X_i} > \delta \right)
     & \leq \frac{1}{n^2 \delta^2} \sum_{i=N + 1}^n \var (X_i)
    \leq \frac{1}{n \delta^2} \sum_{i=N + 1}^n \frac{\var X_i}{i}
    \leq \frac{(n - N) \epsilon}{n \delta^2} \leq \frac{\epsilon}{\delta^2}.
  \end{align*}
  Since $\epsilon > 0$ is arbitrary, this completes the proof.

\end{proof}

\begin{ex}[Durrett 2.2.2]
  The $L^2$ weak law generalizes immediately to certain dependent
  sequences. Suppose $E X_n$ and $E X_n X_m \leq r (n - m)$ for $m \leq n$ (no
  absolute value on the left-hand side!) with $r (k) \to 0$ as $k \to \infty$.
  Show that $(X_1 + \dots + X_n) / n \to 0$ in probability.
\end{ex}

\begin{proof}
  Consider
  \begin{align*}
    E \left( n^{-1} \left( X_1 + \dots + X_n \right) \right)^2
     & = n^{-2} E \left( X_1 + \dots + X_n \right)^2                      \\
     & = n^{-2} \left( \sumi^n E X_i^2 + 2 \sum_{i < j} E X_i E_j \right) \\
     & \leq n^{-2} \left( n r(0) + 2 \sumi^{n - 1} (n - i) r(i) \right).
  \end{align*}
  Let $\epsilon$ be arbitrary.
  Since $r(k) \to 0$ as $k \to \infty$, there is $N$ such that
  $r(k) \leq \epsilon$ for all $k \geq N$. Hence for $n \geq N$, we have
  \begin{align*}
    E \left( n^{-1} (X_1 + \dots + X_n) \right)^2
     & \leq \frac{r(0)}{n} + \frac{2}{n^2} \sumi^N (n - i) r(i)
    + \frac{2}{n^2} \sum_{i = N + 1}^n (n - i) \epsilon         \\
     & \leq \frac{r(0)}{n} + \frac{2 N}{n} \sumi^N r(i)
    + 2 \epsilon.
  \end{align*}
  This shows that $(X_1 + \dots + X_n) / n \to 0$ in $L^2$ and thus
  in probability.

\end{proof}

\begin{ex}[Durrett 2.2.5]
  Let $X_1, X_2, \dots$ be i.i.d.\ with
  \[
    P(X_i > x) = \frac{e}{x \log x}.
  \]
  for $x \geq e$.
  Show that $E \abs{X_i} = \infty$, but there is a sequence of constants
  $\mu_n \to \infty$ so that $S_n / n - \mu_n \to 0$ in probability.
\end{ex}

\begin{proof}
  Note that $P(X_i < e) = 0$. It follows that
  \[
    E \abs{X_i}
    = \int_e^\infty P(X_i > x) \d x
      = [\log \log x]_e^\infty
    = \infty.
  \]
  However,
  \[
    x P(X_i > x) = \frac{e}{\log x} \to 0.
  \]
  Hence, by a previous theorem, if $\mu_n = E X_1 \ind_{(\abs{X_1} \leq n)}$
  then $S_n / n - \mu_n \to 0$. Note that $\mu_n < \infty$ but
  $\mu_n \to E X_i = \infty$ as $n \to \infty$
  by monotone convergence theorem.

\end{proof}

\subsection{Borel-Cantelli lemmas}

\begin{ex}[Durrett 2.3.8]
  Let $\seqinfn{A_n}$ be a sequence of independent events with $P(A_n) < 1$
  for all $n$. Show that $P (\cupinfn A_n) = 1$
  implies $\suminfn P(A_n) = \infty$ and hence $P(A_n \text{ i.o.}) = 1$.
\end{ex}

\begin{proof}
  Suppose for contradiction that $\suminfn P(A_n) = M < \infty$.
  From independence, we have
  \begin{align*}
    P \left( \capn^N A_n^c \right)
    = \prod_{n=1}^N (1 - P(A_n)).
  \end{align*}
  As $\suminfn P(A_n)$ converges, $P(A_n) \to 0$ as $n \to 0$.
  It follows that there exists $\alpha < 1$
  such that $P(A_n) \leq \alpha < 1$ for all $n$.
  Once we select $\beta > 0$ such that $1 - \alpha \geq \exp (\beta \alpha)$,
  we obtain
  \begin{align*}
    P \left( \capn^N A_n^c \right)
    = \prod_{n=1}^N (1 - P(A_n))
    \geq \exp \left( - \beta \sumn^N P(A_n) \right)
  \end{align*}
  This implies
  \begin{align*}
    P \left( \cupinfn A_n \right)
    = 1 - P \left( \capn^\infty A_n^c \right)
    \leq 1 - \exp \left( - \beta \sumn^\infty P(A_n) \right)
  \end{align*}
  Hence, $P (\cupinfn A_n) \leq 1 - \exp (- \beta M) < 1$, contradiction.

\end{proof}

\begin{ex}[Durrett 2.3.9]
  Show that if $P (A_n) \to 0$ and $\suminfn P (A_n^c \cap A_{n + 1}) < \infty$,
  then $P (A_n \text{ i.o.}) = 0$. Find an example of a sequence
  $\seqinfn{A_n}$ to which the result above can
  be applied but the Borel-Cantelli lemma cannot.
\end{ex}

\begin{proof}
  For any $N, M$, we have
  \[
    \bigcup_{n = N}^M A_n
    = A_N \cup \bigcup_{n = N}^{M - 1} (A_n^c \cap A_{n + 1})
  \]
  This implies
  \begin{align*}
    P \left( \bigcup_{n = N}^M A_n \right)
    \leq P (A_N) + \sum_{n = N}^{M - 1} P (A_n^c \cap A_{n + 1}).
  \end{align*}
  Letting $M \to \infty$ shows
  $P (\bigcup_{n = N}^\infty A_n) \leq P(A_n) + \sum_{n = N}^\infty
    P (A_n^c \cap A_{n + 1})$.
  Now letting $N \to \infty$ shows $P (\limsup A_n) = 0$, as desired.

  For an example, define $(\Omega, \fcal, P)$ by letting $\Omega = [0, 1]$,
  $\fcal$ be the Borel sets, and $P$ be the Lebesgue measure.
  For $n \geq 1$, define $A_n$ via
  \[
    A_n = \left[ 1 - 2^{-n} - n^{-1} , 1 - 2^{-n} \right].
  \]
  Then $P(A_n) = n^{-1}$ but $P(A_n^c \cap A_{n + 1}) = 2^{-n - 1}$.


\end{proof}

\begin{ex}[Durrett 2.3.12]
  Let $X_1, X_2, \dots$ be a sequence of r.v.'s on $(\Omega, \fcal, P)$
  where $\Omega$ is a countable set and $\fcal$ consists of all subsets
  of $\Omega$. Show that $X_n \to X$ in probability implies
  $X_n \to X$ a.s.
\end{ex}

\begin{proof}
  Suppose for contradiction that $X_n$ does not converge to $X$ a.s.
  Then there is a set $A \subset \Omega$ such that $P(A) > 0$ and
  $X_n (\omega)$ does not converge to $X(\omega)$ for all $\omega \in A$.
  Let
  \[
    B_r = \left\{ \omega : \abs{X_n (\omega) - X (\omega)}
    > 2^{-r} \text{ i.o.} \right\},
  \]
  then $A = \bigcup_{r=1}^\infty B_r$. By continuity of measure,
  there exists $r$ such that $P (B_r) = \alpha > 0$
  Since $B_r$ is countable, enumerate the elements in $B_r$ by
  $\seqinfi{\omega_i}$. Write $\epsilon = 2^{-r}$.

  Now, there is a subsequence $n^1 (j)$ such that
  $\abs{X_{n^1 (j)} (\omega_1) - X(\omega_1)} > \epsilon$ for all $j$.
  Once we construct $n^{k - 1} (j)$, there is a subsequence $n^{k} (j)$
  such that $\abs{X_{n^k(j)} (\omega_k) - X (\omega_k)} > \epsilon$
  for all $j$. Let $m(k) = n^k (k)$, then
  $\abs{X_{m(k)} (\omega) - X(\omega)} > \epsilon$ for all $k$ and
  all $\omega \in A$. This implies
  \[
    P \left( \abs{X_{m(k)} - X(\omega)} > \epsilon \right)
    = \alpha > 0.
  \]
  Hence, $P (\abs{X_n - X} > \epsilon) = \alpha$ i.o.\ and $X_n$ does not
  converge to $X$ in probability, contradiction.

\end{proof}

\begin{ex}[Durrett 2.3.14]
  Let $X_1, X_2, \dots$ be independent. Show that $\sup X_n < \infty$ a.s.\
  if and only if $\suminfn P(X_n > A) < \infty$ for some $A$.
\end{ex}

\begin{proof}
  ($\implies$) Suppose for contradiction that
  $\suminfn P(X_n > A) = \infty$ for any $A < \infty$.
  By independence and the second Borel-Cantelli lemma,
  we know $P(X_n > A \text{ i.o.}) = 1$ for any $A$.
  Since $A$ is arbitrary, this then implies $\sup X_n = \infty$ a.s.,
  contradiction.

  ($\impliedby$) Suppose $\suminfn P(X_n > A) < \infty$ for some $A$.
  Then, Borel-Cantelli lemma implies $P (X_n > A \text{ i.o.}) = 0$.
  That is, for almost sure $\omega \in \Omega$, there exists
  $N(\omega)$ such that $X_n(\omega) \leq A$ for all $n \geq N$.
  For these $\omega$, $\sup X_n (\omega) < \infty$. Therefore,
  $\sup X_n < \infty$ a.s.

\end{proof}

\section{Central Limit Theorems}

\setcounter{subsection}{2}

\subsection{Characteristic functions}

\begin{ex}[Durrett 3.3.10]
  Using the identity $\sin t = 2 \sin \frac{t}{2} \cos \frac{t}{2}$
  repeatedly leads to
  \[
    \frac{\sin t}{t} = \prod_{m=1}^\infty \cos \left( \frac{t}{2^m} \right).
  \]
  Prove the last identity by interpreting each
  side as a characteristic function.
\end{ex}

\begin{proof}
  Note that
  \[
    \frac{1}{2} \int_{-1}^1 e^{i t x} \d x = \frac{\sin t}{t}.
  \]
  Therefore, LHS can be interpreted as the ch.f. of a
  $\unif [-1, 1]$ random variable. Meanwhile,
  \[
    \cos \left( 2^{-m} t \right) = \frac{1}{2} \exp (i t 2^{-m})
    + \frac{1}{2} \exp (- i t 2^{-m}).
  \]
  Hence, $\cos (2^{-m} t)$ can be interpreted as ch.f. of a random variable
  $X_m$ that maps to $\pm 2^{-m}$ with probability $1/2$ each,
  and RHS can be interpreted as the sum of independent $X_m$'s.
  Now we just need to show that $X_1 + \dots + X_n \convd \unif[-1, 1]$.
  To verify this, note that $X_m = 2^{-m} (2 b_m - 1)$ where
  $\seqinfm{b_m}$ are i.i.d.\ $\Bernoulli (1 / 2)$ random
  variables. Meanwhile, $\suminfm 2^{-m} b_m$ converges to
  $\unif[0, 1]$ by interpreting the sum as choosing binary digits.
  It then immediately follows that $X_1 + \dots + X_n \convd \unif [-1, 1]$.
  More concretely, we can compute the probability of each dyadic intervals
  and approximate all intervals through them. This then uniquely determines
  the distribution.
  By Levy's theorem, the proof for the desired inequality is complete.

\end{proof}

\begin{ex}[Durrett 3.3.20]
  If $X_1, X_2, \dots$ are independent and have characteristic
  function $\exp (- \abs{t}^\alpha)$ then
  $(X_1 + \dots + X_n) / n^{1 / \alpha}$ has the same distribution
  as $X_1$.
\end{ex}

\begin{proof}
  By independence, the ch.f. of $X_1 + \dots + X_n$ is
  $\exp ( - n \abs{t}^\alpha )$. Hence the ch.f. of
  $(X_1 + \dots + X_n) / n^{1 / \alpha}$ is
  \[
    \exp \left( - n \abs{\frac{t}{n^{1 / \alpha}}}^\alpha \right)
    = \exp (- \abs{t}^\alpha ).
  \]
  This is the same ch.f. as $X_1$, so
  $(X_1 + \dots + X_n) / n^{1 / \alpha}$ has the same distribution
  as $X_1$.

\end{proof}

\begin{ex}[Durrett 3.3.23]
  Show that if $X$ and $Y$ are independent and $X + Y$ and
  $X$ have the same distribution then $Y = 0$ a.s.
\end{ex}

\begin{proof}
  Since $X$ and $Y$ are independent, $\phi_{X + Y} = \phi_X \phi_Y$.
  We are given that $\phi_X \phi_Y = \phi_X$, so
  \[
    \phi_X (t) (\phi_Y (t) - 1) = 0
  \]
  for all $t \in \R$. Note that $\phi_X(0) = 1$ and $\phi_X$ is
  continuous, so $\phi_Y = 0$ on a neighborhood of zero.
  Now show that this implies $Y = 0$ a.s.

  Let $t$ be in this neighborhood, then $1 = E[\cos (t Y)]$.
  Note that $\cos (t Y) \leq 1$, so this implies
  $\cos (t Y) = 1$ a.s., and $t Y \in 2 \pi \Z$ a.s. Let
  $t_1, t_2$ be both in this neighborhood and $t_1 / t_2 \notin \Q$.
  Then $t_1 Y \in 2 \pi \Z$ and $t_2 Y \in 2 \pi \Z$ a.s. If
  $Y(\omega) \neq 0$, then
  \[
    \frac{t_1}{t_2} = \frac{t_1 Y(\omega)}{t_2 Y(\omega)} \in \Q,
  \]
  a contradiction. Therefore, $Y = 0$ a.s.

\end{proof}

\section{Martingales}

\subsection{Conditional expectation}

\begin{ex}[Durrett 4.1.2]
  Prove \textbf{Chebyshev's inequality}. If $a > 0$ then
  \[
    P \{ \abs{X} \geq a | \fcal \} \leq a^{-2} E[X^2 | \fcal].
  \]
\end{ex}

\begin{proof}
  Let $A \in \fcal$, we have
  \begin{align*}
    \int_A E[ \ind_{\{\abs{X} \geq a\}} | \fcal] \d P
     & = \int_A \ind_{\{\abs{X} \geq a\}} \d P.
  \end{align*}
  On the event $\{\abs{X} \geq a\}$, $a^{-2} X^2 \geq 1$, so
  \begin{align*}
    \int_A E[ \ind_{\{\abs{X} \geq a\}} | \fcal] \d P
     & \leq \int_A a^{-2} X^2 \d P           \\
     & = a^{-2} \int_A X^2 \d P              \\
     & = a^{-2} \int_A E [X^2 | \fcal] \d P.
  \end{align*}
  Since $A \in \fcal$ is arbitrary, this completes the proof.

\end{proof}

\begin{ex}[Durrett 4.1.5]
  Give an example on $\Omega = \{a, b, c\}$ in which
  \[
    E [E [X | \fcal_1] | \fcal_2] \neq E [E [X | \fcal_2] | \fcal_1].
  \]
\end{ex}

\begin{proof}
  Let $P$ be the uniform distribution and $X = \ind_{\{a\}}$.
  Also let
  \[
    \fcal_1 = \{ \{a, b\}, \{c\}, \emptyset, \Omega\}, \quad
    \fcal_2 = \{ \{b, c\}, \{a\}, \emptyset, \Omega\}.
  \]
  Then $E[X | \fcal_2] = X$. Meanwhile, $E[X | \fcal_1] (c) = 0$
  and $E[X | \fcal_1] (a) = E[X | \fcal_1] (b) = \frac{1}{2}$. With this,
  we can also compute $E[E[X | \fcal_1] | \fcal_2] (a) = \frac{1}{2}$
  and $E[E[X | \fcal_1] | \fcal_2] (b) = E[E[X | \fcal_1] | \fcal_2] (c)
    = \frac{1}{4}$. This shows that $E [E [X | \fcal_1] | \fcal_2] \neq
    E [E [X | \fcal_2] | \fcal_1]$, as desired.

\end{proof}

\begin{ex}[Durrett 4.1.6]
  Show that if $\gcal \subset \fcal$ and $E X^2 < \infty$ then
  \[
    E \left[ (X - E[X | \fcal])^2 \right]
    + E \left[ (E [X | \fcal] - E[X | \gcal])^2 \right]
    = E \left[ (X - E [X | \gcal])^2 \right].
  \]
\end{ex}

\begin{proof}
  For LHS, we have
  \begin{align*}
     & (X - E[X | \fcal])^2 + (E[X | \fcal] - E [X | \gcal])^2      \\
     & = X^2 + 2 E[X | \fcal]^2 - 2 X E[X | \fcal] - 2 E[X | \fcal]
    E [X | \gcal] + E[X | \gcal]^2
  \end{align*}
  For RHS, we have
  \begin{align*}
    (X - E[X | \gcal])^2 = X^2 + E[X | \gcal]^2 - 2 X E[X | \gcal].
  \end{align*}
  Hence, we just need to show
  \[
    E[E[X | \fcal]^2 - X E[X | \fcal] - E[X | \fcal]
        E [X | \gcal]]
    = E[- X E[X | \gcal]].
  \]
  Using the tower property and taking conditional expectation
  w.r.t.\ $\fcal$ first and then taking the unconditional
  expectation, we have
  \begin{align*}
     & E[E[X | \fcal]^2 - X E[X | \fcal] - E[X | \fcal]
           E [X | \gcal]] \\
     & = E[E[X | \fcal]^2] - E[E[X | \fcal]^2]
    - E[E[X | \fcal] E[X | \gcal]]                     \\
     & = - E[E[X | \fcal] E[X | \gcal]],
  \end{align*}
  Similarly, we have
  \begin{align*}
    E[- X E[X | \gcal]]
     & = - E[E[X | \fcal] E[X | \gcal]],
  \end{align*}
  where we used the fact that $E[X | \gcal]$ is $\gcal$-measurable
  and thus $\fcal$-measurable.

\end{proof}

\begin{ex}[Durrett 4.1.8]
  Let $Y_1, Y_2, \dots$ be i.i.d.\ with mean $\mu$ and variance
  $\sigma^2$, $N$ an independent positive integer valued r.v.\
  with $E N^2 < \infty$ and $X = Y_1 + \dots + Y_N$. Show that
  \[
    \var (X) = \sigma^2 E N + \mu^2 \var (N).
  \]
  To understand and
  help remember the formula, think about the two special cases
  in which $N$ or $Y$ is constant.
\end{ex}

\begin{proof}
  It is easy to see that $E[Y_1^2]= \sigma^2 + \mu^2$. Then,
  we have
  \begin{align*}
    E[X] = E[E[X | N]] = E[\mu N] = \mu EN,
  \end{align*}
  and
  \begin{align*}
    E[X^2] = E[E[X^2 | N]]
    = E[N (\sigma^2 + \mu^2) + N (N - 1) \mu^2]
    = \mu^2 E[N^2] + \sigma^2 E N.
  \end{align*}
  Subtrating these two expressions gives $\var X = \sigma^2 EN + \mu^2
    \var N$, as desired.
\end{proof}

\begin{ex}[Durrett 4.1.9]
  Show that if $X$ and $Y$ are random variables with $E[Y | \gcal] = X$
  and $E Y^2 = E X^2 < \infty$, then $X = Y$ a.s.
\end{ex}

\begin{proof}
  Since $X = E[Y | \gcal]$, $X$ is $\gcal$-measurable. It follows that
  \begin{align*}
    E[(Y - X) X | \gcal]
    = E[(Y - X) X | \gcal]
    = X (E[Y | \gcal] - E[X])
    = X (X - X) = 0,
  \end{align*}
  so $E[(Y - X) X] = 0$ i.e.\ $E[XY] = E[X^2]$. Now we have
  \[
    E[(Y - X)^2] = E[Y^2] - 2 E[XY] + E[X^2]
    = E[Y^2] - E[X^2] = 0.
  \]
  Therefore, $X = Y$ a.s.
\end{proof}

\subsection{Martingales, almost sure convergence}

\begin{ex}[Durret 4.2.1]
  Suppose $X_n$ is a martingale w.r.t.\ $\gcal_n$ and let
  $\fcal_n = \sigma (X_1, \dots, X_n)$. Then $\gcal_n \supset
    \fcal_n$ and $X_n$ is a martingale w.r.t\ $\fcal_n$.
\end{ex}

\begin{proof}
  By assumption, $X_1, \dots, X_n$ are all $\gcal_n$ measurable.
  Therefore it is clear that $\gcal_n \supset \fcal_n$.
  Moreover,
  \[
    E[X_{n + 1} | \fcal_n] = E[E[X_{n+1} | \gcal_n] | \fcal_n]
    = E[X_n | \fcal_n] = X_n,
  \]
  so $X_n$ is a martingale w.r.t.\ $\fcal_n$.

\end{proof}

\begin{ex}[Durret 4.2.5]
  Give an example of a martingale $X_n$ with $X_n \to - \infty$ a.s.
  Hint: Let $X_n = \xi_1 + \dots + \xi_n$ where the $\xi_i$ are
  independent (but not identically distributed) with $E \xi_i = 0$.
\end{ex}

\begin{proof}
  Define $\xi_n$ via $P \{\xi_n = - n^{-1}\} = 1 - n^{-2}$
  and $P\{\xi_n = n - n^{-1}\} = n^{-2}$. Then
  \[
    E \xi_n = - \left( 1 - \frac{1}{n^2} \right) \frac{1}{n}
    + \frac{1}{n^2} \left( n - \frac{1}{n} \right) = 0.
  \]
  Therefore $X_n = \xi_1 + \dots + \xi_n$ is a martingale. To show
  $X_n \to -\infty$ a.s., note that
  \[
    \suminfi P\{\xi_n = n - n^{-1}\} = \suminfi \frac{1}{n^2} < \infty.
  \]
  Borel-Cantelli lemma then implies that $P \{\xi_n = n - n^{-1}
    \text{ i.o.}\} = 0$. On the complementary event $\{\xi_n = -n^{-1}
    \text{ for all but finitely many $n$}\}$, it is clear that
  $X_n \to -\infty$.

\end{proof}

\begin{ex}[Durret 4.2.8]
  Let $X_n$ and $Y_n$ be positive integrable and adapted to $\fcal_n$.
  Suppose
  \[
    E[X_{n + 1} | \fcal_n] \leq (1 + Y_n) X_n
  \]
  with $\sum Y_n < \infty$ a.s. Prove that $X_n$ converges a.s. to a finite
  limit by finding a closely related supermartingale to which
  Theorem 4.2.12 can be applied.
\end{ex}

\begin{proof}
  Define $A_n = \prod_{k=1}^{n-1} (1 + Y_k)$ and $Z_n = \frac{X_n}{A_n}$.
  Since $Y_k$ is $\fcal_k$-measurable for each $k$, $A_n$ is
  $\fcal_n$-measurable.
  Also, since $A_n \geq 1$ a.s.\ and $X_n$
  is integrable, $Z_n$ is also integrable. Now we have
  $A_{n + 1} = A_n (1 + Y_n)$, and thus
  \begin{align*}
    E[Z_{n + 1} | \fcal_n]
    = E \left[ \frac{X_{n + 1}}{A_n (1 + Y_n)} \middle| \fcal_n \right]
    = \frac{E[X_{n + 1} | \fcal_n]}{A_n (1 + Y_n)}
    \leq \frac{X_n}{A_n}
    = Z_n.
  \end{align*}
  Therefore, $Z_n$ is a supermartingale and thus converges to some
  finite limit $Z$ a.s.

  Now consider $A_n$. Using $\log (1 + x) \leq x$ for $x > 0$, we have
  \[
    \sumk^{n-1} \log (1 + Y_k)
    \leq \sumk^{n - 1} Y_n
    \leq \suminfn Y_n < \infty \text{ a.s.}
  \]
  Therefore, the series $\sum \log(1 + Y_k)$ converges a.s.
  and thus $A_n = \exp (\sumk^{n-1} \log(1 + Y_k))$ converges
  to a finite limit $A$ a.s. This implies $X_n \to Z A$ a.s.

\end{proof}

\begin{ex}[Durret 4.2.9]
  \textbf{The switching principle}. Suppose $X_n^1$ and $X_n^2$ are
  supermartingales with respect to $\fcal_n$ and $N$ is stopping time
  so that $X_N^1 \geq X_N^2$. Then
  \begin{align*}
    Y_n = X_n^1 \ind_{\{N > n\}} + X_n^2 \ind_{\{N \leq n\}}
    \text{ is a supermartingale.} \\
    Z_n = X_n^1 \ind_{\{N \geq n\}} + X_n^2 \ind_{\{N < n\}}
    \text{ is a supermartingale.} \\
  \end{align*}
\end{ex}

\begin{proof}
  To show $Y_n$ is a supermartingale, note that
  \[
    Y_{n + 1} = X_{n + 1}^1 \ind_{\{N > n + 1\}} + X_{n + 1}^2 \ind_{\{N \leq n + 1\}}
  \]
  However, $\ind_{\{N > n + 1\}} + \ind_{\{N = n + 1\}} = \ind_{\{N > n\}}$
  and $\ind_{\{N \leq n + 1\}} = \ind_{\{N \leq n\}} + \ind_{\{N = n + 1\}}$.
  It follows that
  \[
    Y_{n + 1} = X_{n + 1}^1 \ind_{\{N > n\}}
    + X_{n + 1}^2 \ind_{\{N \leq n\}}
    + (X_{n+1}^2 - X_{n+1}^1) \ind_{\{N = n + 1\}}.
  \]
  Note that $\ind_{\{N \leq n\}}$ and
  $\ind_{\{N > n\}}$ are both $\fcal_n$-measurable, and
  $X_{n+1}^1 \geq X_{n+1}^2$ on the event $\{N = n + 1\}$.
  Taking the conditional expectation w.r.t.\ $\fcal_n$ gives
  \begin{align*}
    E[Y_{n + 1}]
     & \leq \ind_{\{N > n\}} E[X_{n+1}^1 | \fcal_n]
    + \ind_{\{N \leq n\}} E[X_{n+1}^2 | \fcal_n]                   \\
     & \leq \ind_{\{N > n\}} X_n^1 + \ind_{\{N \leq n\}} X_{n+1}^2 \\
     & = Y_n,
  \end{align*}
  showing $Y_n$ is a supermartingale.

  To show $Z_n$ is a supermartingale, note that $\ind_{\{N \geq n + 1\}}$
  and $\ind_{\{N < n + 1\}}$ are $\fcal_n$-measurable. It follows that
  \begin{align*}
    E[Z_{n + 1} | \fcal_n]
     & = E[X_{n + 1}^1 \ind_{\{N \geq n + 1\}} + X_{n + 1}^2 \ind_{\{N < n + 1\}} | \fcal_n]               \\
     & = E[X_{n + 1}^1 | \fcal_n] \ind_{\{N \geq n + 1\}} + E[X_{n + 1}^2  | \fcal_n] \ind_{\{N < n + 1\}} \\
     & \geq X_n^1 \ind_{\{N \geq n + 1\}} + X_n^2 \ind_{\{N < n + 1\}}.
  \end{align*}
  Now we can use a similar decomposition to see
  $E[Z_{n + 1} | \fcal_n] \geq Z_n$.

\end{proof}



\end{document}